% Generisano: uređuj beleske/sNNN.tex i naslove u manifestu.
\documentclass[11pt,a4paper]{article}
% Zajednički preambulum za dokument za čitanje; nije Beamer tema.
\usepackage[a4paper,margin=20mm,headheight=15pt]{geometry}
\usepackage{lmodern,fontspec}
\usepackage[serbian]{babel}
\setmainfont{Latin Modern Sans}
\setsansfont{Latin Modern Sans}
\setmonofont{Latin Modern Mono}
\usepackage{amsmath,amssymb,mathtools,graphicx,xcolor}
\usepackage{array,booktabs,tabularx,multirow,makecell,listings,fancyhdr}
\pagestyle{fancy}\fancyhf{}
\newcommand{\BeleskeTekuciID}{NEOZNACEN}
\fancyhead[L]{\small Beleške uz slajd \texttt{\expandafter\detokenize\expandafter{\BeleskeTekuciID}}}
\fancyfoot[R]{\small\thepage}
\setlength{\parindent}{0pt}\setlength{\parskip}{6pt}
\newwrite\BeleskeMapa
\immediate\openout\BeleskeMapa=\jobname.notes-map.tsv
\AddToHook{shipout/before}{%
  \immediate\write\BeleskeMapa{\BeleskeTekuciID\space\thepage}}
\AddToHook{enddocument/afterlastpage}{\immediate\closeout\BeleskeMapa}
\newcommand{\BeleskeID}[1]{%
  \def\BeleskeZadatiID{#1}%
  \ifx\BeleskeTekuciID\BeleskeZadatiID\else
    \PackageError{beleske}{Pogresan ID beleske: #1}{Proveri mapiranje slajda.}%
  \fi}
\newcommand{\BeleskaSlajda}[4]{%
  \clearpage\gdef\BeleskeTekuciID{#1}%
  {\Large\bfseries Slajd \textnormal{\texttt{\detokenize{#1}}}: #3\par}\medskip
  \includegraphics[page=#2,width=\linewidth]{\BeleskePrezentacija}\par\medskip
  \input{#4}}

\newcommand{\BeleskePrezentacija}{build/08_elementi_analize_logickih_kola.pdf}
\begin{document}
\BeleskaSlajda{08-s001}{1}{Elementi analize logičkih kola}{beleske/s001.tex}
\BeleskaSlajda{08-s002}{2}{Apstraktna logička kola}{beleske/s002.tex}
\BeleskaSlajda{08-s003}{3}{Realna logička kola}{beleske/s003.tex}
\BeleskaSlajda{08-s004}{4}{Statička i dinamička analiza}{beleske/s004.tex}
\BeleskaSlajda{08-s005}{5}{Napajanje}{beleske/s005.tex}
\BeleskaSlajda{08-s006}{6}{Masa}{beleske/s006.tex}
\BeleskaSlajda{08-s007}{7}{Kondenzatori za blokadu}{beleske/s007.tex}
\BeleskaSlajda{08-s008}{8}{Lokalno rasprezanje napajanja}{beleske/s008.tex}
\BeleskaSlajda{08-s009}{9}{Model realnog logičkog kola}{beleske/s009.tex}
\BeleskaSlajda{08-s010}{10}{Statičke karakteristike}{beleske/s010.tex}
\BeleskaSlajda{08-s011}{11}{Raspodela naponskih nivoa}{beleske/s011.tex}
\BeleskaSlajda{08-s012}{12}{Granični ulazni naponi}{beleske/s012.tex}
\BeleskaSlajda{08-s013}{13}{Održanje naponskih nivoa}{beleske/s013.tex}
\BeleskaSlajda{08-s014}{14}{Robusnost}{beleske/s014.tex}
\BeleskaSlajda{08-s015}{15}{Karakteristika prenosa}{beleske/s015.tex}
\BeleskaSlajda{08-s016}{16}{Invertujući i neinvertujući prenos}{beleske/s016.tex}
\BeleskaSlajda{08-s017}{17}{Tri oblasti karakteristike}{beleske/s017.tex}
\BeleskaSlajda{08-s018}{18}{Nelinearne karakteristike}{beleske/s018.tex}
\BeleskaSlajda{08-s019}{19}{Karakteristične tačke}{beleske/s019.tex}
\BeleskaSlajda{08-s020}{20}{Tačka M}{beleske/s020.tex}
\BeleskaSlajda{08-s021}{21}{Regenerabilnost}{beleske/s021.tex}
\BeleskaSlajda{08-s022}{22}{Margine za jednostruki izvor šuma}{beleske/s022.tex}
\BeleskaSlajda{08-s023}{23}{Linearizacija karakteristike}{beleske/s023.tex}
\BeleskaSlajda{08-s024}{24}{Margine za višestruke izvore šuma}{beleske/s024.tex}
\BeleskaSlajda{08-s025}{25}{Pogrešno tumačenje regeneracije}{beleske/s025.tex}
\BeleskaSlajda{08-s026}{26}{Strujni kapacitet}{beleske/s026.tex}
\BeleskaSlajda{08-s027}{27}{Ulazne i izlazne struje}{beleske/s027.tex}
\BeleskaSlajda{08-s028}{28}{Kataloške ulazne struje}{beleske/s028.tex}
\BeleskaSlajda{08-s029}{29}{Opterećenje izlaza}{beleske/s029.tex}
\BeleskaSlajda{08-s030}{30}{Ograničenje izlaznih struja}{beleske/s030.tex}
\BeleskaSlajda{08-s031}{31}{Kataloški izlazni naponi}{beleske/s031.tex}
\BeleskaSlajda{08-s032}{32}{Faktor grananja}{beleske/s032.tex}
\BeleskaSlajda{08-s033}{33}{Logički trud}{beleske/s033.tex}
\BeleskaSlajda{08-s034}{34}{Idealno logičko kolo}{beleske/s034.tex}
\BeleskaSlajda{08-s035}{35}{Vreme uspona i pada}{beleske/s035.tex}
\BeleskaSlajda{08-s036}{36}{Period i kašnjenje}{beleske/s036.tex}
\BeleskaSlajda{08-s037}{37}{Srednje i maksimalno kašnjenje}{beleske/s037.tex}
\BeleskaSlajda{08-s038}{38}{Uticaj ulazne ivice na kašnjenje}{beleske/s038.tex}
\BeleskaSlajda{08-s039}{39}{Uticaj ulazne ivice na izlaznu}{beleske/s039.tex}
\BeleskaSlajda{08-s040}{40}{Idealna pobuda}{beleske/s040.tex}
\BeleskaSlajda{08-s041}{41}{RC kolo}{beleske/s041.tex}
\BeleskaSlajda{08-s042}{42}{Kontinuitet napona kondenzatora}{beleske/s042.tex}
\BeleskaSlajda{08-s043}{43}{Stacionarno stanje RC kola}{beleske/s043.tex}
\BeleskaSlajda{08-s044}{44}{Jednačina RC kola}{beleske/s044.tex}
\BeleskaSlajda{08-s045}{45}{Vremenska konstanta}{beleske/s045.tex}
\BeleskaSlajda{08-s046}{46}{Tevenenov ekvivalent}{beleske/s046.tex}
\BeleskaSlajda{08-s047}{47}{Integracione konstante}{beleske/s047.tex}
\BeleskaSlajda{08-s048}{48}{RC odziv na skok}{beleske/s048.tex}
\BeleskaSlajda{08-s049}{49}{Oblik RC odziva}{beleske/s049.tex}
\BeleskaSlajda{08-s050}{50}{Trajanje RC prelaza}{beleske/s050.tex}
\BeleskaSlajda{08-s051}{51}{Opšte rešenje i početni uslov}{beleske/s051.tex}
\BeleskaSlajda{08-s052}{52}{Opšte rešenje i krajnji uslov}{beleske/s052.tex}
\BeleskaSlajda{08-s053}{53}{Postupak rešavanja RC kola}{beleske/s053.tex}
\BeleskaSlajda{08-s054}{54}{Integrator}{beleske/s054.tex}
\BeleskaSlajda{08-s055}{55}{Diferencijator: početni uslov}{beleske/s055.tex}
\BeleskaSlajda{08-s056}{56}{Diferencijator: odziv}{beleske/s056.tex}
\BeleskaSlajda{08-s057}{57}{Kapacitivno opterećenje}{beleske/s057.tex}
\BeleskaSlajda{08-s058}{58}{Pogrešan proračun impulsa}{beleske/s058.tex}
\BeleskaSlajda{08-s059}{59}{Superpozicija impulsa}{beleske/s059.tex}
\BeleskaSlajda{08-s060}{60}{Promena početnog trenutka}{beleske/s060.tex}
\BeleskaSlajda{08-s061}{61}{Odziv po intervalima}{beleske/s061.tex}
\BeleskaSlajda{08-s062}{62}{Različita trajanja impulsa}{beleske/s062.tex}
\BeleskaSlajda{08-s063}{63}{Diferencijator: impulsni odzivi}{beleske/s063.tex}
\BeleskaSlajda{08-s064}{64}{Induktivnost}{beleske/s064.tex}
\BeleskaSlajda{08-s065}{65}{Postupak rešavanja RL kola}{beleske/s065.tex}
\BeleskaSlajda{08-s066}{66}{RL integrator}{beleske/s066.tex}
\BeleskaSlajda{08-s067}{67}{Kompenzovani razdelnik}{beleske/s067.tex}
\end{document}
